\chapter{The 15-Puzzle}
\index{15 Puzzle}
\section{What it is}
The \emph{15-puzzle} is a sliding puzzle played on a \(4\times 4\) grid. Fifteen
square tiles are labeled \(1\) through \(15\); the sixteenth cell is empty. The
goal is to slide tiles until the numbers appear in order from left to right and
top to bottom, with the empty cell in the bottom-right corner.

\section{Rules}
\begin{itemize}
  \item At any time, exactly one cell is empty.
  \item You may move a tile \emph{into} the empty cell only if that tile shares
  an edge with the empty cell (up, down, left, or right). Diagonal moves are not
  allowed.
  \item Equivalently: you ``slide'' a tile one step orthogonally into the gap;
  the gap moves to where that tile was.
  \item The usual objective is the \emph{solved} configuration: first row
  \(1\)--\(4\), second row \(5\)--\(8\), third row \(9\)--\(12\), fourth row
  \(13\)--\(15\), with the empty cell last.
\end{itemize}

\section{One configuration}
The table below shows one possible arrangement of the board (here the empty
cell is in the bottom-right corner, as in the solved state).

\begin{center}
\begin{tabular}{|c|c|c|c|}
\hline
1 & 2 & 3 & 4 \\
\hline
5 & 6 & 7 & 8 \\
\hline
9 & 10 & 11 & 12 \\
\hline
13 & 14 & 15 & \\
\hline
\end{tabular}
\end{center}


If you want to try the puzzle in a browser, play online at
\url{https://15puzzle.uk/}.

\section{Solving the Puzzle}
\subsection{Small case: \(2\times 2-1\) puzzle}
To see how the \(n\times n\) sliding puzzle behaves, it helps to ask a smaller
question first: not ``what is the shortest move sequence from \(A\) to \(B\)?''
but ``which positions can ever be reached from which?'' The full \(4\times 4\)
game is large; the smallest nontrivial board is \(n=2\), i.e.\ a \(2\times 2\)
grid with three numbered tiles and one blank.

There are \(4! = 24\) ways to place the tiles and the blank. These placements
split into two disjoint \emph{families}: within each family, every configuration
can be turned into every other by some sequence of legal slides, but no legal
sequence connects a position in one family to a position in the other.

\paragraph{One family (``Group~1'').}
Figure~\ref{fig:two-puzzle-group-one} lists all \(12\) configurations in the
family that contains the usual solved board \(1,2,3\) with the blank
bottom-right. 
\begin{figure}[ht]
\centering
\small
\setlength{\tabcolsep}{4pt}
\begin{tabular}{@{}cccc@{}}
\begin{tabular}{|c|c|}
\hline
 & 1 \\
\hline
3 & 2 \\
\hline
\end{tabular}
&
\begin{tabular}{|c|c|}
\hline
 & 2 \\
\hline
1 & 3 \\
\hline
\end{tabular}
&
\begin{tabular}{|c|c|}
\hline
 & 3 \\
\hline
2 & 1 \\
\hline
\end{tabular}
&
\begin{tabular}{|c|c|}
\hline
2 & \\
\hline
1 & 3 \\
\hline
\end{tabular}
\\[1.2em]
\begin{tabular}{|c|c|}
\hline
1 & \\
\hline
3 & 2 \\
\hline
\end{tabular}
&
\begin{tabular}{|c|c|}
\hline
3 & \\
\hline
2 & 1 \\
\hline
\end{tabular}
&
\begin{tabular}{|c|c|}
\hline
1 & 2 \\
\hline
 & 3 \\
\hline
\end{tabular}
&
\begin{tabular}{|c|c|}
\hline
2 & 3 \\
\hline
 & 1 \\
\hline
\end{tabular}
\\[1.2em]
\begin{tabular}{|c|c|}
\hline
3 & 1 \\
\hline
 & 2 \\
\hline
\end{tabular}
&
\begin{tabular}{|c|c|}
\hline
1 & 2 \\
\hline
3 & \\
\hline
\end{tabular}
&
\begin{tabular}{|c|c|}
\hline
2 & 3 \\
\hline
1 & \\
\hline
\end{tabular}
&
\begin{tabular}{|c|c|}
\hline
3 & 1 \\
\hline
2 & \\
\hline
\end{tabular}
\end{tabular}
\caption{All \(12\) configurations in Group~1 (same sliding component as
\(1,2,3\) with blank bottom-right).}
\label{fig:two-puzzle-group-one}
\end{figure}

\paragraph{The other family (``Group~2'').}
Figure~\ref{fig:two-puzzle-group-two} lists the other \(12\) configurations.
No sequence of legal moves joins a board in Figure~\ref{fig:two-puzzle-group-one}
to a board here.

\begin{figure}[ht]
\centering
\small
\setlength{\tabcolsep}{4pt}
\begin{tabular}{@{}cccc@{}}
\begin{tabular}{|c|c|}
\hline
 & 1 \\
\hline
2 & 3 \\
\hline
\end{tabular}
&
\begin{tabular}{|c|c|}
\hline
 & 2 \\
\hline
3 & 1 \\
\hline
\end{tabular}
&
\begin{tabular}{|c|c|}
\hline
 & 3 \\
\hline
1 & 2 \\
\hline
\end{tabular}
&
\begin{tabular}{|c|c|}
\hline
1 & \\
\hline
2 & 3 \\
\hline
\end{tabular}
\\[1.2em]
\begin{tabular}{|c|c|}
\hline
3 & \\
\hline
1 & 2 \\
\hline
\end{tabular}
&
\begin{tabular}{|c|c|}
\hline
2 & \\
\hline
3 & 1 \\
\hline
\end{tabular}
&
\begin{tabular}{|c|c|}
\hline
1 & 3 \\
\hline
 & 2 \\
\hline
\end{tabular}
&
\begin{tabular}{|c|c|}
\hline
2 & 1 \\
\hline
 & 3 \\
\hline
\end{tabular}
\\[1.2em]
\begin{tabular}{|c|c|}
\hline
3 & 2 \\
\hline
 & 1 \\
\hline
\end{tabular}
&
\begin{tabular}{|c|c|}
\hline
1 & 3 \\
\hline
2 & \\
\hline
\end{tabular}
&
\begin{tabular}{|c|c|}
\hline
2 & 1 \\
\hline
3 & \\
\hline
\end{tabular}
&
\begin{tabular}{|c|c|}
\hline
3 & 2 \\
\hline
1 & \\
\hline
\end{tabular}
\end{tabular}
\caption{All \(12\) configurations in Group~2.}
\label{fig:two-puzzle-group-two}
\end{figure}

\subsection{General case: \(n^2-1\) puzzle}
\index{Recursion}
Here a recursive approach is useful. We assume we can
solve the \((n-1)^2-1\) puzzle on the bottom right; we only need
to arrange the numbers on the top row and leftmost column correctly.

By symmetry it suffices to describe how to finish the \emph{first} row; the
first column is analogous.

The first row has \(n\) cells. We can place the first \(n-2\) cells without 
problem. 

The last two cells is tricky. We place the number \(n\) in cell \((n-1)\) 
and make sure the number \(n-1\) is not in cell \(n\). 

\begin{center}
    \small
    \setlength{\tabcolsep}{5pt}
    \renewcommand{\arraystretch}{1.15}
    \begin{tabular}{|c|c|c|c|}
    \hline
    \(1\) & ... & Place n in cell (n-1) & Some other number \\
    \hline
    \(\cdots\) & \(n-1\) is some where here & \(\cdots\) & \(\cdots\) \\
    \hline
    \(\cdots\) & \(\cdots\) & \(\cdots\) & \(\cdots\) \\
    \hline
    \end{tabular}
\end{center}

Next we place the number \(n-1\) in the cell below cell \(n-1\) with number n.
We make the cell \(n\) empty. 
\begin{center}
    \small
    \setlength{\tabcolsep}{5pt}
    \renewcommand{\arraystretch}{1.15}
    \begin{tabular}{|c|c|c|c|}
    \hline
    \(1\) & ... & Place n in cell (n-1) & Some other number \\
    \hline
    \(\cdots\) & & Place \(n-1\) below cell n & \(\cdots\) \\
    \hline
    \(\cdots\) & \(\cdots\) & \(\cdots\) & \(\cdots\) \\
    \hline
    \end{tabular}
\end{center}

Finally we empty the top-right cell~\((1,n)\) as in the last diagram, ready for
the closing moves.
\begin{center}
    \small
    \setlength{\tabcolsep}{5pt}
    \renewcommand{\arraystretch}{1.15}
    \begin{tabular}{|c|c|c|c|}
    \hline
    \(1\) & ... & Place n in cell (n-1) & Make the cell empty \\
    \hline
    \(\cdots\) & & Place \(n-1\) below cell n & \(\cdots\) \\
    \hline
    \(\cdots\) & \(\cdots\) & \(\cdots\) & \(\cdots\) \\
    \hline
    \end{tabular}
\end{center}

From that situation one completes the first row by sliding tile~\(n\) one step
to the right and tile~\(n-1\) one step upward, without disturbing the tiles
\(1,\ldots,n-2\) already fixed at the left end of the row.

The first column is handled in the same way. Applying the recursion to the
remaining lower-right \((n-1)\times(n-1)\) block reduces the problem step by
step until only a \(2\times 2\) puzzle with three tiles remains---the case
already treated above.